\section{Administrasi Basis Data di Produksi: Backup, Restore, dan Pemantauan}

Backup database adalah \textbf{asuransi data} --- tanpa backup, kehilangan data bersifat permanen. Di produksi, backup harus dilakukan secara \textbf{otomatis dan rutin} (minimal harian) menggunakan \code{mysqldump} yang dijadwalkan dengan \code{cron}. Simpan backup di lokasi terpisah (cloud storage, server lain) untuk protection against hardware failure. Uji restore secara berkala --- backup tidak berguna jika file-nya rusak atau prosedur restore tidak teruji \cite{mysqlDocs}.

\begin{contoh}
\textbf{Studi Kasus: Backup Otomatis dan Pemulihan Database}

\textbf{Skenario}: Database toko online perlu di-backup setiap hari jam 02:00 dan disimpan 7 hari terakhir. Suatu hari, data produk terhapus tidak sengaja dan perlu dipulihkan dari backup terakhir.

\textbf{Solusi}: Buat skrip backup yang dijalankan oleh cron; saat terjadi kecelakaan, restore dari file backup.
\end{contoh}

\begin{webcode}[language=PHP, caption={Skrip backup dan cron job}]
#!/bin/bash
# File: /home/deployer/backup-db.sh

DATE=$(date +%Y%m%d_%H%M)
DB_NAME="toko_online"
BACKUP_DIR="/home/deployer/backups"
RETENTION=7

# Backup
mysqldump -u app_user -p"$DB_PASS" $DB_NAME \
  | gzip > "$BACKUP_DIR/${DB_NAME}_${DATE}.sql.gz"

# Hapus backup lebih dari 7 hari
find $BACKUP_DIR -name "*.sql.gz" -mtime +$RETENTION -delete

echo "Backup selesai: ${DB_NAME}_${DATE}.sql.gz"

# Cron: jalankan tiap hari jam 02:00
# crontab -e
# 0 2 * * * /home/deployer/backup-db.sh >> /var/log/backup.log 2>&1

# Restore:
# gunzip < backups/toko_online_20260304.sql.gz | mysql -u root -p toko_online
\end{webcode}

